\documentclass[pdftex,11pt,a4paper] {article}
\usepackage{amsmath,amsfonts,amsthm,amssymb,amscd}
\usepackage{bbm}
\usepackage{latexsym}
\usepackage[pdftex]{graphicx}
%\usepackage[titletoc]{appendix}
\usepackage{hyperref}
%\usepackage[usenames,dvipsnames]{color}

%\usepackage{endnotes,refcount}
\renewcommand{\baselinestretch}{1.0}%0.88%0.905
%\usepackage{makeidx}
%\usepackage[bookmarks=true]{hyperref}
%\usepackage[usenames,dvipsnames]{color}
\usepackage[margin=3.0cm]{geometry}
\setlength{\textheight}{23cm}
\setlength{\textwidth}{16.3cm}
\setlength{\leftmargin}{2.7cm}
\setlength{\rightmargin}{0cm}
\setlength{\topmargin}{-1cm}
%\addtolength{\leftmargin}{0.0cm}


\newtheorem{corollary}{Corollary}[section]
\newtheorem{lemma}{Lemma}[section]
\newtheorem{theorem}{Theorem}[section]
\newtheorem{proposition}{Proposition}[section]
\newtheorem{definition}{Definition}[section]
\newtheorem{notations}{Notations}[section]
\newtheorem{notation}{Notation}[section]
\newtheorem{assumption}{Assumption}[section]
\newtheorem{remark}{Remark}[section]
%\numberwithin{equation}{section}
%\numberwithin{figure}{section}
%\numberwithin{table}{section}
%\parindent=0cm
%\usepackage[font=small,labelfont={small}]{caption}
%\usepackage[hang,small,bf]{caption}



%\usepackage[nottoc,notlot,notlof]{tocbibind}


%\let\footnote=\footnote


\renewcommand{\thefootnote} {\arabic{footnote}}
% \renewcommand\contentsname{}


%\pdfoptionpdfminorversion 6



\begin{document}




\textbf{Conditional Valuation Formula}




Using Lipton (2001), section 13.3, we express the value of a path-dependent option as an integral in the functional space of price trajectories:
\begin{equation} \label {eq:p1}
\begin{split}
& P(0, S_{0}) = e^{-rT} \int_{\Omega}\mathcal{F}(\omega)d\mathcal{D}(\omega)
\end{split}
\end{equation}
where $\mathcal{F}(\omega)$ is a functional mapping trajectories into payoffs and $\mathcal{D}(\omega)$ is the risk-neutral probability measure in the space of trajectories.

We use the augmentation principle as defined in section 13.2 we introduce the functional $J_{t}$  to represent the path-dependent variable linked to the evolution of the spot price $S_{t}$. We consider evaluation of a derivatives security with the terminal pay-off $f(S,J)$.


To model the spot price dynamics, we apply the stochastic volatility model with the instantaneous variance $V_{t}$. Accordingly, we the joint dynamics can be written as follows:
\begin{equation} \label {eq:p2}
\begin{split}
& dS_{t} = rS_{t}dt + \sqrt{V_{t}} S_{t}  \left(\rho dB_{t} + \sqrt{1-\rho^{2}}dW_{t}\right)\\
& dV_{t} = \Phi(V_{t})dt + \Psi(V_{t})dB_{t} \\
& dJ_{0,t} = g_{0}(S_{t}, J_{0,t})dt + g_{1}(S_{t}, J_{0,t})dS_{t}
\end{split}
\end{equation}
where $f$ and $g$ are specified by the augmentation principle.

In this paper we focus on path-dependent barrier options where the conditional on  the lower barrier $B$, $B<S_{0}$:
\begin{equation} \label {eq:p3}
\begin{split}
& J_{t} = \min_{0<t'\leq T} \left( S_{t'} \right), \ dJ_{t} =  \theta \left(J_{t}-S_{t}\right) \left(dS_{t}\right)_{-}, \ f(S,J) = \mathbbm{1}_{ \{ J_{T}> B  \}}f\left(S_{T}\right)
\end{split}
\end{equation}
where $f(S)$ is the terminal payoff function.

The joint dynamics in Eq (\ref{eq:p2}) is Markovian. Accordingly, we can represent the general formula (\ref{eq:p1}) by:
\begin{equation} \label {eq:p4}
\begin{split}
& P(0, S_{0}) = e^{-rT} \int^{\infty}_{V=0}\int^{\infty}_{S=0}\int^{\infty}_{J=0}f(S,J)G(T,V, S,J; 0, V_{0},S_{0}, J_{0})dJdSdV
\end{split}
\end{equation}
where  and $G(T, V S,J; 0, V_{0}, S_{0}, J_{0}) $ is the risk-neutral probability density function of the joint evolution of state variables in SDE (\ref{eq:p2}).


We introduce a novel method to compute the three-dimensional integral in (\ref{eq:p3}).
Conditionally on the filtration generated by Brownian $B_{t}$, $\mathbb{F}^{B}$, we represent the price process by
\begin{equation} \label {eq:p5}
\begin{split}
 & S_{t} = S_{0}\exp\left\{\mu t + q\int^{t}_{0}V_{s}ds + \frac{\rho}{\epsilon}\left(V_{t}- \upsilon\right) + \sqrt{1-\rho^{2}}\int^{t}_{0}\sqrt{V_{s}}dW_{s}\right\}
\end{split}
\end{equation}
Therefore the conditional density of $S_{t}$ is log-normal Gaussian with time-dependent drift and volatility. Accordingly, we represent formula (\ref{eq:p4}) as follows:
 \begin{equation} \label {eq:p4}
\begin{split}
& P(0, S_{0}) = e^{-rT} \mathbb{E}_{V} \left[\int^{\infty}_{S=0}\int^{\infty}_{J=0}f(S,J)G^{LG} \left(T, S,J; 0, S_{0}, J_{0}\left.\right|V(\omega)\right)dJdS \right]
\end{split}
\end{equation}
where the expectation is with respect to the paths of variance process $V_{t}$, and $G^{LG}$ is the log-normal Gaussian density of $S_{t}$ and $J_{t}$ conditional of the variance path $V(\omega)$.

It is clear that for path-independent derivatives, such as European call and put options, the expression in the square brackets simplifies to the Black-Scholes-Merton formula with the conditional variance $V(\omega)$.

We apply MHP to compute the inner integral in closed-form. Thus we generalize Willard's formula for path-dependent options and apply it for valuation of barrier options in closed-form.


\begin{thebibliography}{10}

\bibitem{bn}  Willard A. G. (1997), Calculating Prices and Sensitivities for Path-Independent Derivatives Securities in Multifactor Models, \textit{The Journal of Derivatives}, \textbf{5(1)}, 45-61,



\end{thebibliography}




\end{document}
